\subsection{非紧区间上正函数的积分}

命题 2. 设 f 是一个 \(\geqslant 0\) 的实规则函数，定义在区间 \(I \subset R\) 上，其端点为 a 与 b（a \textless{} b）。为使积分 \(\int_{a}^{b} f(t) dt\) 存在，必须且只需：当 J 跑遍含于 I 的紧区间所成之集时，数集 \(\int_{J} f(t) dt\) 有上界；此时积分 \(\int_{a}^{b} f(t) dt\) 就是数集 \(\int_{J} f(t) dt\) 的上确界。

事实上，由于 \(f \geqslant 0\)，关系 \(J \subset J'\) 蕴含

\[
\int_ {J} f d t \leqslant \int_ {J ^ {\prime}} f d t;
\]

故映射 \(J \mapsto \int_{J} f dt\) 是递增的，而本命题由单调极限定理（Gen.~Top., IV, p.~349, 定理 2）得出。

当映射 \(J \mapsto \int_{J} f(t) dt\) 无界时，它有极限 \(+\infty\)（关于有向集 \(\mathcal{K}(I)\)）；这时人们（滥用语言的说法）说积分 \(\int_{a}^{b} f(t) dt\) 等于 \(+\infty\)。在 \(n^{\circ} 1\) 中建立的积分的性质（当所涉及的是 \(\geqslant 0\) 的函数时）可以推广到所考虑的某些积分为无穷的情形，只要它们所在的关系式有意义。

命题 3（比较原理）。设 f 与 g 是区间 \(I \subset R\) 上的两个实规则函数，使得在 f 与 g 连续的每一点处 \(0 \leqslant f(x) \leqslant g(x)\)（参见 II, p.~61, 命题 6）。若 g 在 I 上的积分收敛，则 f 的积分也收敛，且有 \(\int_{I} f(t) dt \leqslant \int_{I} g(t) dt\)。此外，除非在 I 中 f 与 g 都连续的每一点处 \(f(x) = g(x)\)，否则这两个积分不可能相等。

现在，对每个紧区间 \(\mathrm{J} \subset \mathrm{I}\) 有

\[
\int_ {J} f (t) d t \leqslant \int_ {J} g (t) d t;
\]

由于 \(\int_{\mathrm{J}} g(t) dt \leqslant \int_{\mathrm{I}} g(t) dt\)，诸积分 \(\int_{\mathrm{J}} f dt\) 有上界，故积分 \(\int_{\mathrm{I}} f(t) dt\) 收敛；而且，取极限便得 \(\int_{\mathrm{I}} f(t) dt \leqslant \int_{\mathrm{I}} g(t) dt\)。再设 \(f(x) < g(x)\) 在一个点 \(x \in \mathrm{I}\) 处成立（在该点 \(f\) 与 \(g\) 连续）；存在一个含于 I、不退化为（单个）点的紧区间 \([c, d]\)，使得 \(x \in [c, d]\)；有 \(\int_{c}^{d} f(t) dt < \int_{c}^{d} g(t) dt\)（II, p.~61, 推论 1），又因由上面所述有 \(\int_{a}^{c} f(t) dt \leqslant \int_{a}^{c} g(t) dt\) 与 \(\int_{d}^{b} f(t) dt \leqslant \int_{d}^{b} g(t) dt\)，逐项相加便得 \(\int_{a}^{b} f(t) dt < \int_{a}^{b} g(t) dt\)。

本命题提供了判定函数 \(f \geqslant 0\) 的积分收敛与否的最常用手段，即：把 \(f\) 与一个较简单的函数 \(g \geqslant 0\) 相比较，后者的积分是否收敛已经是知道的；在第五章中我们将看到，在最常见的情形中如何寻找比较函数；并由此外推出积分与级数收敛的常用判别法。

